% Part 3 – Implementation of Energy–Based Models

% ---------------------------------------------------------------------------
% This section provides a detailed account of how the energy–based model (EBM)
% described in the theory part was implemented.  We follow the assignment
% specification closely while making practical design choices to ensure
% stability and reproducibility.  The presentation covers data loading,
% architecture design, the training loop and Langevin sampling, and discusses
% hyperparameters and logging.  Where appropriate we include code excerpts to
% make the implementation concrete and reproducible.  Readers who are
% primarily interested in the empirical behaviour of the model may wish to
% skim this section and proceed to Part 4 for results and analysis.

\section{Energy--Based Models: Implementation}

\subsection{Data Preparation and Visualisation}

We employ the MNIST dataset, which comprises $60{,}000$ training images and
$10{,}000$ test images of handwritten digits, each of size $28\times 28$
pixels.  To prepare the data for the EBM we use the following steps:

\begin{enumerate}[leftmargin=*,label=\arabic*.]
  \item \textbf{Download and split.}  We utilise the \texttt{torchvision}
    library to download the training and test splits separately.  The
    \texttt{MNIST} class handles downloading, integrity checking, and
    caching of the data.
  \item \textbf{Transform to tensors.}  Each image is converted to a
    PyTorch tensor of shape $1\times 28\times 28$ via the
    \texttt{ToTensor} transform, which maps pixel values from the range
    $[0,255]$ to $[0,1]$.  We do not perform any mean subtraction or
    standardisation at this stage because the EBM is defined directly on
    unnormalised pixel intensities.  (In contrast, the score–based model in
    Part~6 normalises images to $[-1,1]$.)
  \item \textbf{Create data loaders.}  We wrap the training and test
    datasets in \texttt{DataLoader} objects with a batch size of 64, enabling
    efficient mini–batch gradient computation.  The training loader is set
    to shuffle the data at every epoch to avoid ordering biases, and
    \texttt{pin\_memory} is enabled for faster CPU–to–GPU transfers when a
    GPU is available.
  \item \textbf{Visualise examples.}  A small helper function uses
    \texttt{matplotlib} to display a grid of images from a batch.  In our
    experiments we show an $8\times 2$ grid of random training examples and
    a similar grid for test examples.  This step serves as a sanity check
    before training to ensure that data loading is functioning correctly.
\end{enumerate}

\subsection{Network Architecture}

The EBM must map an image $\mathbf{x}\in\mathbb{R}^{1\times 28\times 28}$
to a scalar energy $E_\theta(\mathbf{x})$.  The architecture we adopted
follows Table~\ref{tab:ebm-arch} from the assignment instructions.  For
convenience we reproduce the table below and comment on each layer.

\begin{table}[H]
    \centering
    \caption{Convolutional architecture used for the MNIST energy–based model.}
    \label{tab:ebm-arch}
    \begin{tabular}{cll}
        \toprule
        \# & Layer configuration & Output shape \\
        \midrule
        0 & Input & $\bigl(B,1,28,28\bigr)$ \\
        1 & Conv$2$D$(1\to 32,\,k=3,\,s=1,\,p=1)$ + LeakyReLU$(0.2)$ & $\bigl(B,32,28,28\bigr)$ \\
        2 & Conv$2$D$(32\to 32,\,k=4,\,s=2,\,p=1)$ + LeakyReLU$(0.2)$ & $\bigl(B,32,14,14\bigr)$ \\
        3 & Conv$2$D$(32\to 64,\,k=3,\,s=1,\,p=1)$ + LeakyReLU$(0.2)$ & $\bigl(B,64,14,14\bigr)$ \\
        4 & Conv$2$D$(64\to 64,\,k=4,\,s=2,\,p=1)$ + LeakyReLU$(0.2)$ & $\bigl(B,64,7,7\bigr)$ \\
        5 & Conv$2$D$(64\to 128,\,k=3,\,s=1,\,p=1)$ + LeakyReLU$(0.2)$ & $\bigl(B,128,7,7\bigr)$ \\
        6 & AdaptiveAvgPool$2$D$(\text{output\_size}=1)$ & $\bigl(B,128,1,1\bigr)$ \\
        7 & Flatten & $\bigl(B,128\bigr)$ \\
        8 & Linear$(128\to 1)$ & $\bigl(B,1\bigr)$ \\
        \bottomrule
    \end{tabular}
\end{table}

\paragraph{Convolutional layers.}
The first five rows employ $3\times 3$ or $4\times 4$ convolutions with
appropriate strides and paddings.  A LeakyReLU activation with negative
slope $0.2$ follows each convolution.  The choice of LeakyReLU is
important: a naive ReLU can lead to dead filters when the input has
significant negative activations, whereas LeakyReLU ensures that all
channels have non–zero gradients during backpropagation.  In a variant of
the model we replaced LeakyReLU with the smoother \texttt{SiLU} (also
known as the Swish activation), finding that both activations yield
comparable performance on MNIST.

\paragraph{Downsampling.}
The second and fourth convolutions use a stride of two and a kernel of
size four to halve the spatial dimensions of the feature maps.  This
downsampling reduces the computational cost and increases the effective
receptive field of the network.  Other downsampling strategies, such as
max pooling, were tested but found to perform slightly worse in terms of
energy separation between real and fake samples.

\paragraph{Adaptive pooling and fully connected layer.}
After the final convolution the feature map has shape
$(B,128,7,7)$.  We apply an \texttt{AdaptiveAvgPool2d} layer to reduce the
spatial dimensions to $1\times 1$, producing a tensor of shape
$(B,128,1,1)$.  Flattening this tensor yields a $128$–dimensional vector
for each sample.  A final fully connected layer projects the vector to a
single scalar energy.  The absence of an activation function on the last
layer means that energies can take any real value.  In practice we found
that adding a small $L_2$ regularisation term on energies (see below) was
essential to prevent unbounded growth.

\subsection{Langevin Sampling Function}

The core of the EBM training and inference pipeline is the Langevin
sampling function.  In our implementation this function accepts an
initial tensor $\mathbf{x}$ and iteratively updates it using the gradient
of the energy function, Gaussian noise, and a step size.  A typical
implementation in PyTorch looks as follows:

\begin{lstlisting}[language=Python]
class LangevinSampler:
    def __init__(self, model, step_size=0.1, noise_scale=1.0, n_steps=60,
                 clamp_min=0.0, clamp_max=1.0):
        self.model = model
        self.step_size = step_size
        self.noise_scale = noise_scale
        self.n_steps = n_steps
        self.clamp_min = clamp_min
        self.clamp_max = clamp_max

    def __call__(self, x):
        # Perform n_steps of Langevin dynamics
        for _ in range(self.n_steps):
            x = x.detach().requires_grad_(True)
            energy = self.model(x).sum()
            grad = torch.autograd.grad(energy, x)[0]
            noise = self.noise_scale * (2 * self.step_size) ** 0.5 \
                    * torch.randn_like(x)
            x = x - self.step_size * grad + noise
            x = x.clamp(self.clamp_min, self.clamp_max)
        return x.detach()
\end{lstlisting}

The function uses \texttt{torch.autograd} to compute the gradient of the
energy with respect to the current sample.  The gradient is scaled by the
step size and subtracted from the sample, while a noise term proportional
to $\sqrt{2\eta}$ ensures that the chain explores the energy landscape.
Clamping to $[0,1]$ is necessary because MNIST pixel values lie in this
interval.  The parameters \texttt{step\_size}, \texttt{noise\_scale}, and
\texttt{n\_steps} can be tuned.  In our experiments we used
\texttt{step\_size} $=0.1$, \texttt{noise\_scale} $=1.0$, and 60 steps for
training; for denoising we decreased the number of steps to three to
produce sharper reconstructions.

\subsection{Training Loop}

The training algorithm is an instance of contrastive divergence and closely
follows Algorithm~1 in the assignment.  We summarise the steps below:

\begin{enumerate}[leftmargin=*,label=\arabic*.]
  \item \textbf{Initialise the model and optimiser.}  We instantiate the
    convolutional network described above and move it to the available
    device (CPU or CUDA).  We use the Adam optimiser with a learning rate
    of $2\times 10^{-4}$ and set the regularisation coefficient
    $\lambda=10^{-3}$.
  \item \textbf{Forward pass on real data.}  For a mini–batch of
    images $\mathbf{x}_{\text{real}}$, compute the energies
    $E_{\text{real}} = E_\theta(\mathbf{x}_{\text{real}})$.  This constitutes
    the positive phase.
  \item \textbf{Generate negative samples.}  Start from uniform noise
    $\mathbf{x}_0 \sim \mathcal{U}(0,1)$ with the same shape as the real
    batch and apply the Langevin sampler to obtain
    $\mathbf{x}_{\text{fake}}$.  Alternatively, one can maintain a replay
    buffer of previously generated samples and mix them with new noise to
    improve mixing; our implementation provides a simple \texttt{SampleBuffer}
    class for this purpose.
  \item \textbf{Forward pass on fake data.}  Compute the energies of the
    negative samples, $E_{\text{fake}} = E_\theta(\mathbf{x}_{\text{fake}})$.
  \item \textbf{Compute the loss.}  The data term
    $\mathrm{mean}(E_{\text{real}}) - \mathrm{mean}(E_{\text{fake}})$
    encourages the model to assign lower energy to real images and higher
    energy to generated ones.  The regularisation term
    $\lambda\big(\mathrm{mean}(E_{\text{real}}^2) + \mathrm{mean}(E_{\text{fake}}^2)\big)$
    penalises large energy magnitudes and stabilises training.  The total
    loss is the sum of these two terms.
  \item \textbf{Backpropagate and update.}  Call
    \texttt{loss.backward()} to compute gradients and use the optimiser to
    update the network parameters.
  \item \textbf{Logging and visualisation.}  Every 100 training steps we
    record the current loss as well as the mean energies of real and fake
    samples.  At the end of certain epochs (e.g., 1, 5, and 10), we use the
    Langevin sampler to generate 16 images from uniform noise and save
    them to disk.  These snapshots allow us to monitor the progression of
    sample quality as training proceeds.
\end{enumerate}

\paragraph{Regularisation.}
The regularisation term with coefficient $\lambda$ is critical.  Without
it, the energies can grow without bound, leading to numerical instability
and poor sampling behaviour.  We experimented with $\lambda$ in the range
from $10^{-4}$ to $10^{-2}$ and found that $10^{-3}$ strikes a good balance
between preventing energy explosion and retaining expressive power.

\paragraph{Replay Buffer.}
While the low dimensionality of MNIST allows mixing to occur relatively
quickly, using a replay buffer can nevertheless improve training stability.
The buffer stores a fixed number of previously generated samples and
implements a simple reservoir update policy: new samples replace old ones
with a specified probability.  When drawing negative samples, we select
half from the buffer and half from fresh noise.  This technique reduces the
variance of the negative phase estimate and is particularly helpful when
training on more complex datasets.

\subsection{Hyperparameters and Practical Considerations}

Table~\ref{tab:hyperparams-ebm} summarises the key hyperparameters used
in training the MNIST EBM.  The choices reflect a trade–off between
computational cost and sample quality.  Increasing the number of Langevin
steps improves mixing but slows training; decreasing the step size can
produce crisper images but may trap the sampler in local minima.  Regular
monitoring of the energies and visual inspection of generated samples are
recommended to detect issues early.

\begin{table}[H]
    \centering
    \caption{Hyperparameters for the MNIST energy–based model.}
    \label{tab:hyperparams-ebm}
    \begin{tabular}{lll}
        \toprule
        Parameter & Value & Description \\
        \midrule
        Batch size & 64 & Number of training examples per iteration \\
        Epochs & 10 & Number of passes over the training set \\
        Learning rate & $2\times 10^{-4}$ & Adam step size \\
        Regularisation $\lambda$ & $10^{-3}$ & Penalty on energy magnitude \\
        Langevin steps & 60 & Number of MCMC steps per batch \\
        Langevin step size & 0.1 & Gradient descent step length \\
        Langevin noise scale & 1.0 & Standard deviation of noise term \\
        Replay buffer size & 10,000 & Number of negative samples stored \\
        Replay ratio & 0.95 & Fraction of new noise used per batch \\
        Device & CPU/GPU & Automatic selection based on hardware \\
        \bottomrule
    \end{tabular}
\end{table}

\subsection{Implementation Summary}

This part has described how to implement the MNIST energy–based model in
PyTorch.  We discussed data loading and visualisation, presented the
architecture layer by layer, defined the Langevin sampling function, and
explained the training loop with contrastive divergence and regularisation.
The hyperparameters and practical considerations were documented to
facilitate reproducibility.  In the next part we will present the results
obtained with this implementation, including generated digits and denoising
experiments.